% =========================================================================
% config.tex -- master configuration for the BOOK layout
% -------------------------------------------------------------------------
% Mirrors config_manuscript.tex: same package set (where relevant) and the
% SAME custom-command NAMES, so a manuscript chapter needs only its call
% PARAMETERS updated (done by transform_manuscript.py), not its structure.
%
% Everything that is "a number you might want to iterate on" is exposed as
% a \newcommand near the top of the relevant section. Change it here once.
% =========================================================================

% -----------------------------------------------------------------------
% Packages  (same family as config_manuscript.tex)
% -----------------------------------------------------------------------
\usepackage[utf8]{inputenc}
\usepackage[T1]{fontenc}
\usepackage{amsmath, amssymb}
\usepackage{graphicx}
\usepackage{booktabs}
\usepackage{threeparttable}
\usepackage{array}
\usepackage[english]{babel}
\usepackage{hyperref}   % must load before cleveref
\usepackage{cleveref}
\usepackage{xstring}
\usepackage{geometry}
\usepackage{mdframed}
\usepackage{tcolorbox}       % styled coloured boxes
  \tcbuselibrary{skins, breakable}
\usepackage{tikz}            % needed for \bookcropimage (bio portrait center-crop)
\usepackage{xcolor}
\newif\ifbooksoulloaded
\IfFileExists{soul.sty}{\usepackage{soul}\booksoulloadedtrue}{\booksoulloadedfalse}
\ifbooksoulloaded\else\providecommand{\hl}[1]{\textcolor{red}{#1}}\fi
\usepackage{setspace}
\usepackage{microtype}
\usepackage[shortlabels]{enumitem}
\usepackage{ragged2e}
% Fonts
\usepackage{helvet}
\usepackage{times}
\usepackage{titlesec}   % section/chapter title formatting
\usepackage{eurosym}
\usepackage{changepage}
\usepackage{wasysym}
\usepackage{fancyhdr}
\usepackage{lastpage}
\usepackage{eso-pic}    % full-bleed cover image
\usepackage{multicol}   % balanced spoiler columns on the chapter cover
\IfFileExists{lipsum.sty}{\usepackage{lipsum}}{} % blind text for mockups
\usepackage{stfloats}   % REQUIRED: allows [b] placement for full-width
                        % (figure*/table*) floats -- standard LaTeX supports
                        % only [t]/[p] for those. The take-home box, the
                        % biography box and L-size bottom figures all need it.

% Editorial line numbering, kept for parity with the manuscript config but
% OFF by default in the book layout. Flip \booklinenumberstrue to turn on.
\usepackage{lineno}
\renewcommand{\linenumberfont}{\normalfont\tiny\sffamily\color{gray}}
\newif\ifbooklinenumbers
\booklinenumbersfalse
\ifbooklinenumbers\AtBeginDocument{\linenumbers}\fi



\newcommand{\degree}{\ensuremath{{}^{\circ}}}
\newcommand{\authornote}[1]{\hl{[Author note: #1]}}

% Set sans-serif font
\renewcommand{\sfdefault}{phv} % phv=Helvetica


% =========================================================================
% 1. PAGE GEOMETRY -- page size, columns, margins. All explicit, editable.
%    (Heading/paragraph/float spacing lives in section 3, not here.)
% =========================================================================

% ---- Page size -----------------------------------------------------------
\newcommand{\bookpagewidth}{216mm}
\newcommand{\bookpageheight}{276mm}

% ---- Two-column text block ------------------------------------------------
\newcommand{\bookcolwidth}{83.5mm}     % single column width
\newcommand{\bookcolsep}{5mm}          % gap between columns
\newcommand{\booktextwidth}{172mm}     % = 2*colwidth + colsep (172 = 2*83.5+5)

% ---- Lateral margins (symmetric; derived from page/text width) -----------
% (\bookpagewidth - \booktextwidth) / 2 = (216-172)/2 = 22mm
\newcommand{\booklatmargin}{22mm}

% ---- Vertical margins ------------------------------------------------------
\newcommand{\booktopblank}{5mm}       % blank strip above the header box
\newcommand{\bookheadheight}{7mm}      % height of the header text box
\newcommand{\bookheadsep}{3mm}         % gap between header rule and body text
% total top margin to first text line = 16+7+2 = 25mm (matches spec)
\newcommand{\booktopmargintotal}{25mm}
\newcommand{\bookbottommargin}{20mm}   % blank strip below body text (no footer)

\setlength{\columnsep}{\bookcolsep}

% `includehead` makes `top` measure from the paper edge to the TOP OF THE
% HEADER, so the header lives INSIDE the top margin (as specified):
%   0 .. 16mm      blank
%   16 .. 23mm     header text box (\bookheadheight)
%   23 .. 25mm     headsep
%   25mm           first body line   (= \booktopmargintotal)
\geometry{
  paperwidth   = \bookpagewidth,
  paperheight  = \bookpageheight,
  left         = \booklatmargin,
  right        = \booklatmargin,
  top          = \booktopblank,
  headheight   = \bookheadheight,
  headsep      = \bookheadsep,
  includehead,
  bottom       = \bookbottommargin,
  footskip     = 1pt,
  twoside,
}

% Text body tab / paragraph indent (same names as manuscript config).
% Paragraph VERTICAL spacing (\parskip) is set in section 3.
\newcommand{\bodytab}{10mm}
\setlength{\parindent}{4mm}
\linespread{1.0}


% =========================================================================
% 2. COLOURS
% =========================================================================
\definecolor{maintextcolor}{RGB}{33,30,31}
\definecolor{headingcolor}{RGB}{42, 110, 184}
\definecolor{boxcolor}{RGB}{244, 247, 252}
\definecolor{bookcovernumcolor}{RGB}{205,205,205}
\definecolor{bookcovercapcolor}{RGB}{111,111,111}

% -----------------------------------------------------------------------
% Part-dependent box colour
% -----------------------------------------------------------------------
% Part 1: Chapters 0-5   -> light blue
% Part 2: Chapters 6-9   -> light green
% Part 3: Chapters 10-15 -> light amber
\definecolor{boxcolorpart1}{RGB}{244, 247, 252}
\definecolor{boxcolorpart2}{RGB}{238, 252, 250}
\definecolor{boxcolorpart3}{RGB}{252, 246, 238}

\newcommand{\boxcolor}{%
  \ifnum\value{chapter}<6
    \colorlet{boxcolorcurrent}{boxcolorpart1}%
  \else\ifnum\value{chapter}<10
    \colorlet{boxcolorcurrent}{boxcolorpart2}%
  \else
    \colorlet{boxcolorcurrent}{boxcolorpart3}%
  \fi\fi
}
\colorlet{boxcolorcurrent}{boxcolorpart1} % safe default before first \boxcolor call

% Header rule (full page width; colour/thickness parameterised)
\colorlet{headerrulecolor}{headingcolor}
\newcommand{\headerrulethickness}{1.5pt}

\usepackage[font={footnotesize,sf}, labelfont={bf,color=headingcolor}, labelsep=quad]{caption}


% =========================================================================
% 3. HEADINGS, PARAGRAPHS & FLOAT SPACING
% -------------------------------------------------------------------------
% SIMPLE APPROACH, ON PURPOSE: every before/after space below is a single
% plain \newcommand you can edit directly -- no measuring, no rounding, no
% automatic snapping. \bookvunit is kept only as a convenient unit to write
% values in (e.g. "1.5\bookvunit" reads as "one and a half body-text
% lines") -- it is NOT read back or used to compute anything automatically.
% If two-column pages drift out of alignment at the bottom, that is a
% separate, later problem to solve downstream (e.g. by hand-tuning the
% values here per chapter, or picking values that happen to be exact
% multiples of \bookvunit) -- this section does not attempt to solve it
% for you automatically.
% =========================================================================

% ---- 3.1  Grid unit (convenience only, not auto-applied) ------------------
\newlength{\bookvunit}
\AtBeginDocument{\setlength{\bookvunit}{\baselineskip}}

% ---- 3.2  Heading spacing: BEFORE and AFTER, all editable here ------------
\newcommand{\bookspacebeforesection}{1.5\bookvunit}
\newcommand{\bookspaceaftersection}{0.5\bookvunit}

\newcommand{\bookspacebeforesubsection}{1.2\bookvunit}
\newcommand{\bookspaceaftersubsection}{0.4\bookvunit}

\newcommand{\bookspacebeforesubsubsection}{1.0\bookvunit}
\newcommand{\bookspaceaftersubsubsection}{0.3\bookvunit}

% Paragraph (run-in heading): before is vertical, after is a horizontal
% gap between the bold label and the text that follows on the same line.
\newcommand{\bookspacebeforeparagraph}{1.0\bookvunit}
\newcommand{\bookspaceafterparagraph}{0.5em}

% ---- 3.3  Paragraph vertical spacing (regular body paragraphs) ------------
\newcommand{\bookparskip}{0.0\bookvunit}
\AtBeginDocument{\setlength{\parskip}{\bookparskip}}

% ---- 3.4  Numbering depth & label formats ---------------------------------
\setcounter{secnumdepth}{3}
\renewcommand{\thesubsubsection}{(\alph{subsubsection})}

% ---- 3.5  Label box widths per level, editable here ------------------------
% \thesection is short ("1", "12"); \thesubsection is wider ("1.2",
% "12.10"); \thesubsubsection is short again ("(a)", "(b)"). Each level
% gets its own box width so the number never crowds the title text.
\newcommand{\booksecnumwidth}{2.2em}
\newcommand{\booksubsecnumwidth}{3.2em}
\newcommand{\booksubsubsecnumwidth}{2.6em}

% ---- 3.6  Chapter (kept minimal/unused) ------------------------------------
% Unlike the manuscript, the BOOK chapter heading is not auto-printed by
% \chapter (see \startchapter in section 5) -- it is hand-composed on the
% chapter's "first page" (\chapterfirstpage, section 7) so it can sit
% beside the chapter-number box.
\titleformat{\chapter}[block]{}{}{0pt}{}

% ---- 3.7  Section ----------------------------------------------------------
\titleformat{\section}[block]
  {\bfseries\fontfamily{phv}\selectfont\large\color{headingcolor}}
  {\makebox[\booksecnumwidth][l]{\thesection}}
  {0pt}
  {}
\titlespacing*{\section}
  {0pt}{\bookspacebeforesection}{\bookspaceaftersection}

% ---- 3.8  Subsection -------------------------------------------------------
\titleformat{\subsection}[block]
  {\bfseries\fontfamily{phv}\selectfont\normalsize\color{headingcolor}}
  {\makebox[\booksubsecnumwidth][l]{\thesubsection}}
  {0pt}
  {}
\titlespacing*{\subsection}
  {0pt}{\bookspacebeforesubsection}{\bookspaceaftersubsection}

% ---- 3.9  Subsubsection -----------------------------------------------------
\titleformat{\subsubsection}[block]
  {\bfseries\fontfamily{phv}\selectfont\small\color{headingcolor}}
  {\makebox[\booksubsubsecnumwidth][l]{\thesubsubsection}}
  {0pt}
  {}
\titlespacing*{\subsubsection}
  {0pt}{\bookspacebeforesubsubsection}{\bookspaceaftersubsubsection}

% ---- 3.10  Paragraph (run-in heading) ---------------------------------------
\titleformat{\paragraph}[runin]
  {\bfseries\fontfamily{phv}\selectfont\small\color{maintextcolor}}
  {}
  {0pt}
  {}
\titlespacing*{\paragraph}
  {0pt}{\bookspacebeforeparagraph}{\bookspaceafterparagraph}

% ---- 3.11  Float spacing (figures, tables, boxes) ---------------------------
% Standard LaTeX registers controlling space AROUND floats. All exposed
% here as named, editable commands rather than being buried in defaults.
%   \bookfloatsep      -- space BETWEEN two floats stacked on the same
%                          page (e.g. two figures both placed at [t]).
%   \booktextfloatsep   -- space BETWEEN a float and the body text next to
%                          it on the same page (e.g. below a [t] figure,
%                          above a [b] figure).
%   \bookintextsep       -- space above/below an in-text (non-float-page)
%                          float when it interrupts the running text.
\newcommand{\bookfloatsep}{0.25\bookvunit}
\newcommand{\booktextfloatsep}{0.5\bookvunit}
\newcommand{\bookintextsep}{1.0\bookvunit}
\AtBeginDocument{%
  \setlength{\floatsep}{\bookfloatsep}%
  \setlength{\textfloatsep}{\booktextfloatsep}%
  \setlength{\intextsep}{\bookintextsep}%
  \setlength{\dblfloatsep}{\bookfloatsep}%
  \setlength{\dbltextfloatsep}{\booktextfloatsep}%
}

% ---- 3.12  Body text color / links ------------------------------------------
\color{maintextcolor}

\hypersetup{
  colorlinks  = true,
  linkcolor   = maintextcolor,
  citecolor   = blue!60!black,
  urlcolor    = blue!60!black,
}


% =========================================================================
% 4. PAGE STYLE / HEADERS  (header only, no footer)
% -------------------------------------------------------------------------
% LEFT (even) page:  "<page num, bold> | Chapter <N>  <chapter short name>"
%                     left-aligned at the outer (left) margin.
% RIGHT (odd) page:  "<section num> <section title> | <page num, bold>"
%                     right-aligned at the outer (right) margin.
% Both fully in \headingcolor.
% =========================================================================

% Page number: bold, headingcolor (used on both sides).
\newcommand{\bookheadpage}[1]{\sffamily\scriptsize\bfseries\color{headingcolor}#1}

% Running section mark: captures BOTH the section number and title so the
% header can format them together ("2.3  Light Reactions") without relying
% on \rightmark's default uppercase behaviour. Fed by \sectionmark below.
\makeatletter
\newcommand{\booksectionmarkstore}{}
\renewcommand{\sectionmark}[1]{%
  \renewcommand{\booksectionmarkstore}{\thesection\ \ #1}%
  \markright{#1}% keep \rightmark populated too, for anything else relying on it
}
\makeatother

\fancypagestyle{bookheader}{%
  \fancyhf{}%
  % ---- Left (even) page: page number | Chapter N  short name ------------
  \fancyhead[LE]{%
    \sffamily\scriptsize\color{headingcolor}%
    \bookheadpage{\thepage}\quad\textbar\quad\mdseries Chapter~\thechapter\ \ \bookchaptershortstore%
  }%
  % ---- Right (odd) page: section num  title | page number ----------------
  \fancyhead[RO]{%
    \sffamily\scriptsize\color{headingcolor}%
    \booksectionmarkstore\quad\textbar\quad\bookheadpage{\thepage}%
  }%
  \renewcommand{\headrulewidth}{\headerrulethickness}%
  \renewcommand{\headrule}{%
    \hspace*{-\booklatmargin}%
    {\color{headerrulecolor}\rule[0pt]{\bookpagewidth}{\headerrulethickness}}%
  }%
  \renewcommand{\footrulewidth}{0pt}%
}
\pagestyle{bookheader}

% Cover pages get no header/rule/page-number at all.
\fancypagestyle{bookcover}{%
  \fancyhf{}%
  \renewcommand{\headrulewidth}{0pt}%
  \renewcommand{\headrule}{}%
  \renewcommand{\footrulewidth}{0pt}%
}

\makeatletter
\newcommand{\bookchaptershortstore}{}
\newcommand{\chaptershortname}[1]{\renewcommand{\bookchaptershortstore}{#1}}
\newcommand{\bookchaptertitlestore}{}
\makeatother


% =========================================================================
% 5. CHAPTER-START-ON-EVEN-PAGE  (parameterised)
% =========================================================================
% Toggle which parity chapters open on. Default: EVEN, per spec.
\newif\ifchapterstartseven
\chapterstartseventrue

\makeatletter
\newcommand{\cleartoevenpage}{%
  \clearpage
  \if@twoside
    \ifodd\c@page
      \hbox{}\thispagestyle{empty}\newpage
      \if@twocolumn\hbox{}\thispagestyle{empty}\newpage\fi
    \fi
  \fi
}
\newcommand{\cleartooddpage}{%
  \clearpage
  \if@twoside
    \ifodd\c@page\else
      \hbox{}\thispagestyle{empty}\newpage
      \if@twocolumn\hbox{}\thispagestyle{empty}\newpage\fi
    \fi
  \fi
}
\newcommand{\clearchapterpage}{\ifchapterstartseven\cleartoevenpage\else\cleartooddpage\fi}
\makeatother

% Registers the chapter (counter, ToC, running marks, resets figure/table/
% equation counters) WITHOUT printing anything -- the visible heading is
% hand-built by \chaptercover / \chapterfirstpage below.
\makeatletter
\newcommand{\startchapter}[1]{%
  \clearchapterpage
  \refstepcounter{chapter}%
  \renewcommand{\bookchaptertitlestore}{#1}%
  \renewcommand{\bookchaptershortstore}{#1}% default; override with \chaptershortname
  \addcontentsline{toc}{chapter}{\protect\numberline{\thechapter}#1}%
  \markboth{#1}{}%
  \boxcolor
  \setcounter{workedexample}{0}%
}
% Preface / unnumbered front matter (still registers a ToC entry + marks,
% but is not counted against the part-colour chapter numbering).
\newcommand{\startchapterstar}[1]{%
  \clearchapterpage
  \renewcommand{\bookchaptertitlestore}{#1}%
  \renewcommand{\bookchaptershortstore}{#1}%
  \addcontentsline{toc}{chapter}{#1}%
  \markboth{#1}{}%
  \boxcolor
}
\makeatother


% =========================================================================
% 6. CHAPTER COVER PAGE  (opening page: full-bleed image + number/caption
%    bar + two-column "in this chapter" spoiler)
% =========================================================================
\newcommand{\bookcoverimageheight}{160mm}   % full-bleed photo height (from page top)
\newcommand{\bookcoverbarheight}{32mm}      % number/caption bar height -- tie to font below
\newcommand{\bookcovernumwidth}{32mm}       % number box width (kept ~= bar height => "square-ish")
\newcommand{\bookcovernumfont}{\fontsize{70}{75}\selectfont}
\newcommand{\bookcovercapwidth}{137mm}
\newcommand{\bookcovergaptobar}{4mm}        % gap: bar -> spoiler container
\newcommand{\bookcoverlabelgap}{3mm}        % gap: "In this chapter..." label -> columns
\newcommand{\bookcoverspoilerheight}{70mm}  % FIXED total height of the spoiler container
\newcommand{\bookcovercapspread}{1.3}

% Full-bleed cover photograph. Tries T1 first, then T0, each across
% pdf -> png -> jpeg -> jpg:
%   ./figures/chapter{NN}/T1/{label}_T1.(pdf|png|jpeg|jpg)
%   ./figures/chapter{NN}/T0/{label}_T0.(pdf|png|jpeg|jpg)
%
% IMPORTANT (fixes two earlier bugs):
%  1. The extracted chapter number is stored via \xdef into
%     \bookcovchapnum -- a GLOBAL, fully-expanded plain-text macro. \xdef
%     (unlike \def) fully expands its replacement text at definition time
%     and assigns globally, so \bookcovchapnum reliably holds the plain
%     digits (e.g. "01") in every scope that reads it afterwards, instead
%     of silently reverting to undefined once a local \StrLeft/\StrRight
%     group closes.
%  2. \bookcovchapnum is used with NO shared "whole path" macro name
%     (no \coverfigbaseT1/\coverfigbaseT0) -- every path is written out in
%     full at each \IfFileExists/\includegraphics call, so there is
%     nothing left for a chapter file to collide with by redefining a
%     same-named macro elsewhere.
\makeatletter
\newcommand{\makecoverimage}[1]{%
  \StrLeft{#1}{3}[\covtmpA]\StrRight{\covtmpA}{2}[\covtmpB]%
  \xdef\bookcovchapnum{\covtmpB}%
  \AddToShipoutPictureBG*{%
  \AtPageUpperLeft{\raisebox{-\bookcoverimageheight}{%
    \IfFileExists{./figures/chapter\bookcovchapnum/T1/#1_T1.pdf}{%
      \includegraphics[width=\bookpagewidth,height=\bookcoverimageheight]{./figures/chapter\bookcovchapnum/T1/#1_T1.pdf}%
    }{\IfFileExists{./figures/chapter\bookcovchapnum/T1/#1_T1.png}{%
      \includegraphics[width=\bookpagewidth,height=\bookcoverimageheight]{./figures/chapter\bookcovchapnum/T1/#1_T1.png}%
    }{\IfFileExists{./figures/chapter\bookcovchapnum/T1/#1_T1.jpeg}{%
      \includegraphics[width=\bookpagewidth,height=\bookcoverimageheight]{./figures/chapter\bookcovchapnum/T1/#1_T1.jpeg}%
    }{\IfFileExists{./figures/chapter\bookcovchapnum/T1/#1_T1.jpg}{%
      \includegraphics[width=\bookpagewidth,height=\bookcoverimageheight]{./figures/chapter\bookcovchapnum/T1/#1_T1.jpg}%
    }{\IfFileExists{./figures/chapter\bookcovchapnum/T0/#1_T0.pdf}{%
      \includegraphics[width=\bookpagewidth,height=\bookcoverimageheight]{./figures/chapter\bookcovchapnum/T0/#1_T0.pdf}%
    }{\IfFileExists{./figures/chapter\bookcovchapnum/T0/#1_T0.png}{%
      \includegraphics[width=\bookpagewidth,height=\bookcoverimageheight]{./figures/chapter\bookcovchapnum/T0/#1_T0.png}%
    }{\IfFileExists{./figures/chapter\bookcovchapnum/T0/#1_T0.jpeg}{%
      \includegraphics[width=\bookpagewidth,height=\bookcoverimageheight]{./figures/chapter\bookcovchapnum/T0/#1_T0.jpeg}%
    }{\IfFileExists{./figures/chapter\bookcovchapnum/T0/#1_T0.jpg}{%
      \includegraphics[width=\bookpagewidth,height=\bookcoverimageheight]{./figures/chapter\bookcovchapnum/T0/#1_T0.jpg}%
    }{%
      \colorbox{gray!25}{\parbox[b][\bookcoverimageheight][c]{\bookpagewidth}{%
        \centering\sffamily\large\color{gray!60}[cover photo: #1]}}%
    }}}}}}}}%
  }}%
}%
  % Reserve the part of the image that intrudes into the normal text block.
  % The text block already starts \booktopmargintotal below the page top, so
  % only the remainder has to be pushed down.
  \rule{0pt}{\dimexpr\bookcoverimageheight-\booktopmargintotal\relax}\par
}
\makeatother

% Number / caption bar directly beneath the cover photo. Occupies exactly
% \bookpagewidth (full page bleed for the background colour) but keeps its
% CONTENT (number + caption) constrained within \booktextwidth, matching
% the rest of the book's text block:
%   [ \booklatmargin | number zone | caption \bookcovercapwidth | \booklatmargin ]
% Caption is right-aligned within the text block, justified, in
% \bookcovercapcolor with no paragraph indent. Number occupies whatever
% text-width space remains to the left of the caption, centred within
% that zone, in \bookcovernumcolor. Both cells share ONE background colour
% (\boxcolorcurrent, the current part colour) so the bar reads as a single
% continuous full-width band with no internal seam.
\newsavebox{\bookcovercapbox}
\newlength{\bookcoverbarht}


\newcommand{\makecoverbar}[1]{%
  \sbox{\bookcovercapbox}{%
    \parbox[c][][c]{\bookcovercapwidth}{%
      \noindent \setlength{\parindent}{0pt}%
      \linespread{\bookcovercapspread}\selectfont
      \sffamily\footnotesize\color{bookcovercapcolor}\justifying #1}}%
  \setlength{\bookcoverbarht}{\dimexpr\ht\bookcovercapbox+\dp\bookcovercapbox+4mm\relax}%
  \ifdim\bookcoverbarht<\bookcoverbarheight
    \setlength{\bookcoverbarht}{\bookcoverbarheight}%
  \fi
  \noindent\hspace*{-\booklatmargin}%
  \colorbox{boxcolorcurrent}{%
    \parbox[c][\bookcoverbarht][c]{\booklatmargin}{\hfill}% left page-margin strip, background only
    \parbox[c][\bookcoverbarht][c]{%
        \dimexpr\booktextwidth-\bookcovercapwidth\relax}{%
      \centering\sffamily\bookcovernumfont\color{bookcovernumcolor}\thechapter}%
    \parbox[c][\bookcoverbarht][c]{\bookcovercapwidth}{%
      \usebox{\bookcovercapbox}}%
    \parbox[c][\bookcoverbarht][c]{\booklatmargin}{\hfill}% right page-margin strip, background only
  }%
  \par
}

% "In this chapter ..." label -- STANDALONE single-column element,
% width-constrained to column 1 (\bookcolwidth), sitting above the
% multicols body instead of being its first line.
\newcommand{\bookcoverlabel}{%
  \parbox{\bookcolwidth}{\noindent\textit{In this chapter} \ldots}%
}

% Spoiler body text -- no longer prints the label itself.
\newcommand{\inthischapter}[1]{%
\noindent #1

}

% Full spoiler container: label, gap, then the two-column body -- all
% inside a fixed-height \bookcoverspoilerheight parbox.
\newcommand{\makecoverspoiler}[1]{%
  \parbox[b][\bookcoverspoilerht][b]{\booktextwidth}{%
    \bookcoverlabel

    \vspace{\bookcoverlabelgap}

    \begin{multicols}{2}
      \inthischapter{#1}%
    \end{multicols}
  }%
}
\newlength{\bookcoverspoilerht}

\newcommand{\chaptercover}[3]{%
  \twocolumn[{%
    \thispagestyle{bookcover}%
    \makecoverimage{#1}%
    \makecoverbar{#2}%
    \vspace{\bookcovergaptobar}%
    \setlength{\bookcoverspoilerht}{\dimexpr
      \bookpageheight
      -\bookbottommargin
      -\bookcoverimageheight
      -\bookcoverbarht
      -\bookcovergaptobar
    \relax}%
    \makecoverspoiler{#3}%
  }]%
  \thispagestyle{bookcover}%
  \clearpage
}


% =========================================================================
% 7. CHAPTER FIRST PAGE  (title + numbered box, top; skills box top of
%    the second column)
% =========================================================================
\newcommand{\bookfirstpagetitlesize}{\fontsize{24}{26}\selectfont}
\newcommand{\bookfirstpagenumbox}{16mm}
\newcommand{\bookfirstpagetitlegap}{10mm}
\newcommand{\bookfirstpagenumboxwidth}{30mm}   % rectangle width
\newcommand{\bookfirstpagenumboxheight}{35mm}  % rectangle height, from the header rule down

% #1 = professional-skills content (an itemize, typically) that will sit
% inside a \boxskills box at the top of the SECOND column.
% The \twocolumn header block now carries ONLY the title row (title left,
% number box right), so column 1 starts immediately beneath it instead of
% being held down by the height of the skills box.
%
% The skills box is an ORDINARY single-column float, pinned to the top of
% the page and to the SECOND column:
%   \suppressfloats[t] clears the top of the CURRENT column only (column 1);
%   \@topnum is restored from \topnumber at the start of every column, so
%   column 2 is still free to take it. Declaring the float immediately after
%   means column 2's top is the first slot available to it.
\newcommand{\chapterfirstpage}[1]{%
  \twocolumn[{%
    \thispagestyle{bookheader}%
    \vspace*{-\dimexpr\bookheadsep-0.15mm\relax}%
    \noindent
    \begin{minipage}[c]{\dimexpr\booktextwidth-\bookfirstpagenumboxwidth-4mm\relax}
      \raggedright\bookfirstpagetitlesize\bfseries\fontfamily{phv}\selectfont
      \color{headingcolor}\bookchaptertitlestore\par
    \end{minipage}%
    \hfill
    \colorbox{boxcolorcurrent}{\parbox[c][\bookfirstpagenumboxheight][c]{\bookfirstpagenumboxwidth}{%
      \centering\sffamily\bookcovernumfont\color{bookcovernumcolor}\thechapter}}%
    \par\vspace{\bookfirstpagetitlegap}%
  }]%
  \thispagestyle{bookheader}%
  \suppressfloats[t]%
  \begin{figure}[t]
    \begin{boxskills}
    #1
    \end{boxskills}
  \end{figure}%
}

% =========================================================================
% 8. ENDING PAGE  /  HOMEWORK PAGES
% =========================================================================
% boxmessages always spans both columns and sits at the bottom of the page
% -- this is INVARIANT per spec, so it is baked into the environment and
% chapter files need no float wrapper at the call site.
\newtcolorbox{boxmessagesinner}{
  myboxstyle, width=\linewidth,
  title = {Take home messages},
}

\newenvironment{boxmessages}
  {\onecolumn
   \thispagestyle{bookheader}%
   \begin{boxmessagesinner}}
  {\end{boxmessagesinner}%
  \vspace*{\fill}%
   \twocolumn}

% Convenience break used before the take-home box and before "Homework"
% (\newpage only breaks columns in twocolumn mode; \clearpage guarantees
% a fresh page, which is what the ending / homework pages need).
\newcommand{\chapterhardbreak}{\clearpage}

% =========================================================================
% 9. FIGURES  (S = one column, L = two columns, M = 126mm + side caption)
% =========================================================================
% ---- Figure image-block SIZES -------------------------------------------
% Width is fixed by the size class; HEIGHT is now also fixed, so a page's
% vertical budget is predictable regardless of the artwork's aspect ratio
% (and so placeholders occupy exactly the space the real image will).
\newcommand{\bookfigwidthS}{\linewidth}   % one column
\newcommand{\bookMfigwidth}{126mm}        % medium, side caption
\newcommand{\bookfigwidthL}{\linewidth}   % two columns (= \booktextwidth)

% Width is fixed by the size class. HEIGHT below is a PLACEHOLDER-ONLY
% spec: it sizes the gray fallback box when no image file is found, so a
% missing image still reserves a realistic footprint. Once a real image
% is found, height is NOT enforced -- \bookfitimage scales it by width
% only and lets its natural aspect ratio set the height.
\newcommand{\bookfigheightS}{70mm}
\newcommand{\bookfigheightM}{100mm}
\newcommand{\bookfigheightL}{120mm}

% ---- Figure caption spacing (image -> caption, and caption -> following
% content). Applies to S and L (caption sits BELOW the image); M has a
% SIDE caption so these skips don't apply there.
\newcommand{\bookfigcapabove}{0.8\bookvunit}   % gap: image bottom -> caption top
\newcommand{\bookfigcapbelow}{0.2\bookvunit}   % gap: caption bottom -> whatever follows


% Images inside boxes / bio portraits keep their own heights.
\newcommand{\bookfiglabelS}{S \textbar{} 83.5 $\times$ 70 mm}
\newcommand{\bookfiglabelM}{M \textbar{} 126 $\times$ 100 mm}
\newcommand{\bookfiglabelL}{L \textbar{} 172 $\times$ 120 mm}

\newcommand{\bookboxfigheight}{45mm}
\newcommand{\bookbiofigheight}{25mm}    % was 40mm -- now the FIXED crop height
\newcommand{\bookbiofigwidth}{18.75mm}     % 3:4 portrait -> width = 0.75 * height
\newcommand{\bookbiotitlegap}{3mm}   % gap: name/dates line -> body text, right pane
\newcommand{\bookbiogap}{2mm}   % gap: image pane -> text pane


% ---- Center-crop ("cover") loader -- fills a FIXED width x height box
% completely, cropping any overflow evenly from the centre, instead of
% letterboxing (\bookfitimage) or leaving height uncontrolled
% (\bookfitimagewidth). Needs the image's OWN natural size to decide
% which dimension to scale by, so it measures it via \sbox first.
\newsavebox{\bookcropimgbox}
\newlength{\bookcropimgnw}
\newlength{\bookcropimgnh}
\newcommand{\bookcropimage}[3]{%
  % #1 = image path  #2 = target width  #3 = target height
  \sbox{\bookcropimgbox}{\includegraphics{#1}}%
  \setlength{\bookcropimgnw}{\wd\bookcropimgbox}%
  \setlength{\bookcropimgnh}{\ht\bookcropimgbox}%
  \begin{tikzpicture}
    \path[use as bounding box] (0,0) rectangle (#2,#3);
    \clip (0,0) rectangle (#2,#3);
    \pgfmathtruncatemacro{\bookcropimgwider}{%
      (#2/#3) > (\bookcropimgnw/\bookcropimgnh) ? 1 : 0}%
    \ifnum\bookcropimgwider=1
      \node[inner sep=0pt] at (0.5*#2,0.5*#3) {\includegraphics[width=#2]{#1}};
    \else
      \node[inner sep=0pt] at (0.5*#2,0.5*#3) {\includegraphics[height=#3]{#1}};
    \fi
  \end{tikzpicture}%
}

% Tier/extension fallback loader for the crop path -- SAME search order
% (T1 -> T0, pdf -> png -> jpeg -> jpg) as \bookincludeimage, but calls
% \bookcropimage on a hit instead of \bookfitimage.
\newcommand{\bookincludeimagecrop}[5]{%
  \bookincludeimagecropbase{#1}{#3}{#4}{#5}{#2}%
}
\newcommand{\bookincludeimagecropbase}[5]{%
  \IfFileExists{#1.pdf}{\bookcropimage{#1.pdf}{#3}{#4}}{%
  \IfFileExists{#1.png}{\bookcropimage{#1.png}{#3}{#4}}{%
  \IfFileExists{#1.jpeg}{\bookcropimage{#1.jpeg}{#3}{#4}}{%
  \IfFileExists{#1.jpg}{\bookcropimage{#1.jpg}{#3}{#4}}{%
    \IfStrEq{#5}{}{%
      \bookimageplaceholder{#2}{#3}{#4}{portrait}% no more bases
    }{%
      \bookincludeimagecropbase{#5}{#2}{#3}{#4}{}% try the other tier
    }%
  }}}}%
}


% ---- Unified image loader: TIER fallback (T1 -> T0) then EXTENSION
% ---- fallback (pdf -> png -> jpeg -> jpg), used by every image command.
% #1 = T1 base path (folder+filename, no extension), e.g.
%      ./figures/chapter02/T1/C02-fig-x_T1
% #2 = T0 base path (folder+filename, no extension), e.g.
%      ./figures/chapter02/T0/C02-fig-x_T0
% #3 = label shown on placeholder   #4 = target width   #5 = target height
% #6 = human-readable size label shown on placeholder
\newcommand{\bookincludeimage}[6]{%
  \bookincludeimagebase{#1}{#3}{#4}{#5}{#6}{#2}%
}
% Tries ONE base across all 4 extensions; #6 = fallback base to try next
% (the OTHER tier), empty {} once there is nothing left to fall back to.
\newcommand{\bookincludeimagebase}[6]{%
  \IfFileExists{#1.pdf}{\bookfitimage{#1.pdf}{#3}{#4}}{%
  \IfFileExists{#1.png}{\bookfitimage{#1.png}{#3}{#4}}{%
  \IfFileExists{#1.jpeg}{\bookfitimage{#1.jpeg}{#3}{#4}}{%
  \IfFileExists{#1.jpg}{\bookfitimage{#1.jpg}{#3}{#4}}{%
    \IfStrEq{#6}{}{%
      \bookimageplaceholder{#2}{#3}{#4}{#5}% no more bases -> placeholder
    }{%
      \bookincludeimagebase{#6}{#2}{#3}{#4}{#5}{}% try the other tier
    }%
  }}}}%
}

% #3 (target height) is kept in the signature only so the call site doesn't
% need to change -- it is now used SOLELY by \bookimageplaceholder (the
% gray fallback box) below. Once a real image file is found, only WIDTH is
% fixed; the image keeps its own natural aspect ratio and the slot's
% height is whatever that produces.
\newcommand{\bookfitimage}[3]{%
  \parbox[t]{#2}{\centering\includegraphics[width=#2,keepaspectratio]{#1}}%
}
\newcommand{\bookimageplaceholder}[4]{%
  \fbox{\parbox[c][\dimexpr#3-2\fboxsep-2\fboxrule\relax][c]{%
        \dimexpr#2-2\fboxsep-2\fboxrule\relax}{%
    \centering\sffamily\footnotesize\color{gray!70}
    Figure\\#1\\[2mm]{\tiny #4}}}%
}

% Main figure command -- SAME NAME as the manuscript, extended with two
% leading optional args: size (S/M/L, default L) and placement (t/b,
% default t). Manuscript calls \figplaceholder{label}{name}{caption} keep
% working unchanged (they just get the L/t default).
% #1=size #2=pos #3=label(also used to find ./figures/chapterNN/T1|T0/#3_T1|T0.*)
% #4=display name shown over the fallback box  #5=caption text
%
% Same fix as \makecoverimage above: the extracted chapter number goes
% into a GLOBAL, fully-expanded macro (\bookfigchapnum) via \xdef, and
% paths are written out in full at each call site -- no shared
% \figbaseT1/\figbaseT0 macro name, so nothing to collide with.
\NewDocumentCommand{\figplaceholder}{O{L} O{t} m m m}{%
  \StrLeft{#3}{3}[\figtmpA]\StrRight{\figtmpA}{2}[\figtmpB]%
  \xdef\bookfigchapnum{\figtmpB}%
  \IfStrEq{#1}{S}{%
    \begin{figure}[#2]
    \centering
    \bookincludeimage
      {./figures/chapter\bookfigchapnum/T1/#3_T1}%
      {./figures/chapter\bookfigchapnum/T0/#3_T0}%
      {#4}{\bookfigwidthS}{\bookfigheightS}{\bookfiglabelS}
    \captionsetup{aboveskip=\bookfigcapabove, belowskip=\bookfigcapbelow}%
    \caption{#5}\label{#3}
    \end{figure}%
  }{%
  \IfStrEq{#1}{M}{%
    \begin{figure*}[#2]
    \noindent
    % Both minipages are [t] AND both open with \vspace{0pt}. The \vspace is
    % what actually does the work: without it, a minipage whose first item is
    % a box (here the fixed-height image block, a \parbox[c]) takes that box's
    % CENTRE as its reference baseline, so the caption would sit level with
    % the middle of the image instead of its top edge.
    \begin{minipage}[t]{\bookMfigwidth}
      \vspace{0pt}%
      \centering\bookincludeimage
        {./figures/chapter\bookfigchapnum/T1/#3_T1}%
        {./figures/chapter\bookfigchapnum/T0/#3_T0}%
        {#4}{\linewidth}{\bookfigheightM}{\bookfiglabelM}
    \end{minipage}\hfill%
    \begin{minipage}[t]{\dimexpr\booktextwidth-\bookMfigwidth-\bookcolsep\relax}
      \vspace{0pt}%
      \captionof{figure}{#5}\label{#3}
    \end{minipage}
    \end{figure*}%
  }{%
    \begin{figure*}[#2]
    \centering
    \bookincludeimage
      {./figures/chapter\bookfigchapnum/T1/#3_T1}%
      {./figures/chapter\bookfigchapnum/T0/#3_T0}%
      {#4}{\bookfigwidthL}{\bookfigheightL}{\bookfiglabelL}
    \captionsetup{aboveskip=\bookfigcapabove, belowskip=\bookfigcapbelow}%
    \caption{#5}\label{#3}
    \end{figure*}%
  }}%
}


% Width-only image loader: fits the image to a target WIDTH only, height
% follows naturally from the image's own aspect ratio (no keepaspectratio
% height cap, no letterboxing). Used where the image should always visibly
% fill the available width, e.g. \boxfigplaceholder.
\newcommand{\bookfitimagewidth}[2]{%
  \parbox{#2}{\centering\includegraphics[width=#2]{#1}}%
}

% Width-only variant of the tier/extension fallback loader -- same search
% order (T1 -> T0, pdf -> png -> jpeg -> jpg) but sizes to WIDTH only.
% #1 = T1 base  #2 = T0 base  #3 = placeholder label  #4 = target width
\newcommand{\bookincludeimagewidth}[4]{%
  \bookincludeimagewidthbase{#1}{#3}{#4}{#2}%
}
\newcommand{\bookincludeimagewidthbase}[4]{%
  \IfFileExists{#1.pdf}{\bookfitimagewidth{#1.pdf}{#3}}{%
  \IfFileExists{#1.png}{\bookfitimagewidth{#1.png}{#3}}{%
  \IfFileExists{#1.jpeg}{\bookfitimagewidth{#1.jpeg}{#3}}{%
  \IfFileExists{#1.jpg}{\bookfitimagewidth{#1.jpg}{#3}}{%
    \IfStrEq{#4}{}{%
      \bookimageplaceholder{#2}{#3}{\bookboxfigheight}{box image}% no more bases
    }{%
      \bookincludeimagewidthbase{#4}{#2}{#3}{}% try the other tier
    }%
  }}}}%
}

\newcommand{\boxfigplaceholder}[1]{%
  \StrLeft{#1}{3}[\figtmpA]\StrRight{\figtmpA}{2}[\figtmpB]%
  \xdef\bookfigchapnum{\figtmpB}%
  \begin{center}
  \bookincludeimagewidth
    {./figures/chapter\bookfigchapnum/T1/#1_T1}%
    {./figures/chapter\bookfigchapnum/T0/#1_T0}%
    {#1}{\linewidth}%
  \end{center}
}

\newcommand{\boxbiofigplaceholder}[1]{%
  \StrLeft{#1}{3}[\figtmpA]\StrRight{\figtmpA}{2}[\figtmpB]%
  \xdef\bookfigchapnum{\figtmpB}%
  \centering
  \bookincludeimagecrop
    {./figures/chapter\bookfigchapnum/T1/#1_T1}%
    {./figures/chapter\bookfigchapnum/T0/#1_T0}%
    {#1}{\bookbiofigwidth}{\bookbiofigheight}%
}

% =========================================================================
% 10. BOXES  (S/L sized via the call site -- see transform_manuscript.py,
%     which wraps \begin{boxmain}...\end{boxmain} in figure/figure*[t/b])
% =========================================================================

% Inner box padding, split lateral vs vertical so they can be tuned
% independently (previously left/right were hardcoded 1mm and top/bottom
% were left to tcolorbox defaults -- this made vertical look inconsistent
% with lateral, since it wasn't actually the same value).
\newcommand{\bookboxpad}{1mm}        % global padding
\newcommand{\bookboxpadlat}{0.5mm}        % left/right, inside the frame
\newcommand{\bookboxpadtop}{0.5mm}        % body content: gap above first line
\newcommand{\bookboxpadbottom}{0.5mm}     % body content: gap below last line
\newcommand{\bookboxpadtoptitle}{0.5mm}   % title bar: gap above title text
\newcommand{\bookboxpadbottomtitle}{0mm}% title bar: gap below title, to its rule

\tcbset{
  myboxstyle/.style={
    enhanced,
    % NOT breakable: every box in the book layout lives inside a float
    % (figure/figure*), and floats cannot break across pages -- a breakable
    % tcolorbox inside one renders corrupted, overlapping the body text.
    % A box that will not fit on one page must be split editorially.
    breakable = false,
    colback      = boxcolorcurrent,
    colbacktitle = boxcolorcurrent,
    coltitle     = headingcolor,
    colframe     = boxcolorcurrent,
    boxrule      = 0pt,
    arc          = 0pt,
    boxsep       = \bookboxpad,
    left         = \bookboxpadlat,
    right       = \bookboxpadlat,
    top          = \bookboxpadtop,
    bottom      = \bookboxpadbottom,
    toptitle     = \bookboxpadtoptitle,
    bottomtitle = \bookboxpadbottomtitle,
    fonttitle    = \bfseries\sffamily\small,
    fontupper    = \sffamily\footnotesize\color{maintextcolor},
    fontlower    = \sffamily\footnotesize\color{maintextcolor},
  }
}

% The actual numbered tcolorbox -- SAME keys/theme/title signature as the
% manuscript's boxmain, just renamed \boxmaininner. \boxmain (below) is
% now the public name and owns the size/placement decision itself.
\newtcolorbox[
  auto counter,
  number within = chapter,
  list inside   = lob,
]{boxmaininner}[3][]{
  myboxstyle, width=\linewidth,
  title = {Box~\thetcbcounter\;#2: #3},
  #1,
}
\crefname{tcb@cnt@boxmaininner}{box}{boxes}
\Crefname{tcb@cnt@boxmaininner}{Box}{Boxes}

% Boolean recording which literal float (figure vs figure*) THIS \boxmain
% call opened, so the matching \end{...} closes the same one. Kept as a
% plain \if...\fi (not \IfStrEq) because \if...\fi does no grouping of
% its own -- \begin{figure}'s internal group can stay open safely all the
% way across to the matching \end{figure} in the end-code below. \IfStrEq
% cannot safely wrap a literal \begin{...}/\end{...} directly (its own
% internal grouping breaks the float's group balance). And unlike
% \expandafter\begin\expandafter{...}, every \begin/\end name below is a
% plain literal string, so the linter can resolve and match all of them.
\newif\ifbookboxsizeS

\NewDocumentEnvironment{boxmain}{O{L} O{t} O{} m m}{%
  \IfStrEq{#1}{S}{\bookboxsizeStrue}{\bookboxsizeSfalse}%
  \ifbookboxsizeS
    \begin{figure}[#2]
  \else
    \begin{figure*}[#2]
  \fi
  \begin{boxmaininner}[#3]{#4}{#5}%
}{%
  \end{boxmaininner}
  \ifbookboxsizeS
    \end{figure}
  \else
    \end{figure*}
  \fi
}

% Professional-skills box (unchanged) -- placed by \chapterfirstpage.
\newtcolorbox{boxskills}{
  myboxstyle,
  title = {Professional practice relevance and skills developed},
}

% Biographies -- ALWAYS two-column width, bottom of page (invariant), so
% the float wrapper is baked in here and chapter files call it exactly as
% in the manuscript: \begin{boxbio}{Name}{birth}{death}...\end{boxbio}
\newtcolorbox{boxbioinner}[3]{
  myboxstyle, width=\linewidth,
  breakable=false, sidebyside, sidebyside align=top,
  lefthand width=\bookbiofigwidth,
  sidebyside gap=\bookbiogap,
  before lower={%
    {\bfseries\sffamily\small\color{headingcolor}Biography: #1 (#2 -- #3)}\par
    \vspace{\bookbiotitlegap}%
  },
}

% Optional leading [t/b] controls page placement, default b per spec
% ("biography always at the bottom of the page"). Existing manuscript
% calls with no bracket (\begin{boxbio}{Name}{birth}{death}) keep working
% unchanged, since the default already matches the old hardcoded [b].
\newenvironment{boxbio}[4][b]
  {\begin{figure*}[#1]\begin{boxbioinner}{#2}{#3}{#4}}
  {\end{boxbioinner}\end{figure*}}

% =========================================================================
% 11. WORKED EXAMPLES  (inline, not floated -- unchanged from manuscript)
% =========================================================================
\newcounter{workedexample}
\numberwithin{workedexample}{chapter}
\renewcommand{\theworkedexample}{\thechapter.\arabic{workedexample}}

\newenvironment{workedexample}[1]
{%
  \refstepcounter{workedexample}%
  \setlength{\parindent}{0pt}%
  \par\medskip\noindent
  {\bfseries\small\fontfamily{phv}\selectfont\color{headingcolor}%
   Worked Example~\theworkedexample: #1}\par
}
{%
  \par\hfill End of Worked Example $\blacksquare$\par
}
\newcommand{\subcase}[1]{\par\noindent\textbf{#1.}\enspace\ignorespaces}
\crefname{workedexample}{Worked Example}{Worked Examples}
\Crefname{workedexample}{Worked Example}{Worked Examples}

% -------------------------------------------------------------
% Per-chapter numbering (unchanged from manuscript)
\numberwithin{equation}{chapter}
\numberwithin{figure}{chapter}
\numberwithin{table}{chapter}


% =========================================================================
% 12. FLOAT PLACEMENT PARAMETERS
% =========================================================================
% This is a figure-dense book in which figures/boxes may only sit at the
% top or bottom of a page, so LaTeX must be allowed to fill much more of a
% page with floats than the (very conservative) defaults permit.
\setcounter{topnumber}{3}
\setcounter{bottomnumber}{2}
\setcounter{totalnumber}{4}
\setcounter{dbltopnumber}{3}
\renewcommand{\topfraction}{0.9}
\renewcommand{\bottomfraction}{0.7}
\renewcommand{\dbltopfraction}{0.9}
\renewcommand{\textfraction}{0.08}
\renewcommand{\floatpagefraction}{0.7}
\renewcommand{\dblfloatpagefraction}{0.7}

